\section{Security and trust}
\label{sec:security}

The spec says \texttt{rd\_data} is only valid while the FIFO is not empty and
does not say what an empty FIFO shows, so the scoreboard never checks it. An
empty FIFO that still shows old data is a leak: data a reset was meant to
discard could reach the next user. We added one coverpoint there (vplan S1).

\paragraph{What we cover and measure.} Every read request at empty, alone or
racing a write, right after a reset and after a drain: four bins, all hit by a
short directed sequence. Each time, the testbench checks whether
\texttt{rd\_data} shows a word written only before the latest reset. Memory
starts unknown, so such a word must be left over.

\paragraph{What we found.} The clean RTL leaks. In the full suite, 2439 of
7058 read requests at empty showed pre-reset data. Reset clears the pointers
and count but not the memory, and \texttt{rd\_data} always shows the slot at
the read pointer. The spec allows this, so the default run reports the count
and passes. With \texttt{FIFO\_TRUST\_STRICT=1} it is an assertion, and the
clean RTL fails at 10{,}011~ns. A hardware fix would drive \texttt{rd\_data}
to zero while empty. The released pool stays 8/8, and the run matches on
Eustis. Logs are in \path{evidence/security/}.

\paragraph{Limits.} Only the read side is checked, not how long old data stays
in memory.
